\chapter{Quick Reference Guide to GNU/Linux Forensic Commands}

This appendix provides a practical quick-reference cheat sheet containing essential command-line tools that forensic examiners and incident responders utilize when auditing Unix and GNU/Linux systems, during both live response triage and static post-mortem disk analysis.

\section{Live System Triage and State Inspection}

\subsection{Running Processes and Volatile State}
\begin{lstlisting}[style=terminal]
# List all running processes with user, PID, resource usage, and full command-line:
ps auxwww

# Hierarchical parent-child process tree (reveals orphan or injected processes):
pstree -p -a

# Open file descriptors and active network sockets for a specific process (e.g., PID 1234):
lsof -p 1234

# Verify actual binary executable on disk invoked by the process:
ls -l /proc/1234/exe

# Extract process environment variables from memory:
strings -a /proc/1234/environ | tr '\0' '\n'

# Dump virtual memory mapping for a process:
cat /proc/1234/maps
\end{lstlisting}

\subsection{Network Sockets and Active Connections}
\begin{lstlisting}[style=terminal]
# Listening and established TCP/UDP sockets with PID and program name:
netstat -tulpn
# Modern alternative (iproute2):
ss -tulpan

# List all processes maintaining active network connections:
lsof -i -n -P

# Raw packet capture on a specific network interface (bypassing DNS resolution):
tcpdump -i eth0 -nn -s 0 -w incident_capture.pcap
\end{lstlisting}

\subsection{Privileged Binaries and Special Attributes}
\begin{lstlisting}[style=terminal]
# Locate all files with SUID or SGID permission bits enabled:
find / -type f \( -perm -04000 -o -perm -02000 \) -exec ls -l {} \; 2>/dev/null

# Search for non-root owned files inside system binary paths:
find /bin /sbin /usr/bin /usr/sbin /dev -not -user root -ls

# Recursive audit of immutable (+i) and append-only (+a) Ext2/3/4 attributes:
lsattr -aR /dev /bin /sbin /usr/lib /etc 2>/dev/null | grep -E "(\-i\-|\-a\-)"
\end{lstlisting}

\section{Forensic Image Analysis with The Sleuth Kit (TSK)}

When examining a raw bit-stream disk or partition image (\texttt{disk.img}):

\begin{lstlisting}[style=terminal]
# 1. General filesystem telemetry (block size, inode range, metadata counts):
fsstat disk.img

# 2. Recursive file listing with associated inodes (including deleted entries):
fls -r -p -m "/" disk.img > bodyfile.txt

# 3. In-depth inspection of a specific inode (MAC timestamps, block pointers):
istat disk.img 533487

# 4. Direct extraction of inode content into an external output file:
icat disk.img 533487 > carved_binary.bin

# 5. Enumerate deleted and orphan inodes only:
ils -r -m disk.img >> bodyfile.txt

# 6. Generate formatted chronological timeline (mactime):
mactime -b bodyfile.txt -d -z CST > forensic_timeline.csv
\end{lstlisting}

\section{Integrity Verification and Anomaly Detection}

\begin{lstlisting}[style=terminal]
# Compute cryptographic hashes for evidentiary integrity preservation:
md5sum disk.img > disk.img.md5
sha256sum disk.img > disk.img.sha256

# Extract printable ASCII and Unicode string constants from a suspicious binary:
strings -a -t x suspicious_binary.bin

# Package integrity audit on Debian/Ubuntu systems:
debsums -c -s

# Package integrity audit on Red Hat/CentOS/Fedora systems:
rpm -Va
\end{lstlisting}
